%% kapa: exam=1 round=19 year=2008 period=2 form=A subject=english subject_ko=영어 seq=40 number=80
%% source: Q-Net 감정평가사 기출문제 (cst00309) artlSeq=1001962 file=extracted/1008/1001962/2008년 제19회 감정평가사 1차시험 기출문제/2교시.hwp
%% conversion: kordoc 4.24.51 -> tools/gen1.py (auto); review=proofread-sonnet; flags=none
\begin{question}{80}
다음 글을 읽고 보기의 문장이 들어갈 위치로 가장 알맞은 곳을 고르시오.
\begin{framedtext}
Every day millions of people type messages on their computers in order to make small talk, discuss current events, fall in love, or find new friends. \textbf{(A)} Each virtual community has people with similar interests, ideas, or questions. \textbf{(B)} One community may focus on a health issue, while another may try to make political changes. \textbf{(C)} Just like neighbors in small towns, virtual community members may start out talking about one topic and end up talking about another. \textbf{(D)} They may also offer help when a community member needs it. \textbf{(E)}
\end{framedtext}
\begin{bogi}[보기]
These people form "virtual communities" on the Internet.
\end{bogi}
\begin{choices}
  \item (A)
  \item (B)
  \item (C)
  \item (D)
  \item (E)
\end{choices}
\end{question}
\answer{1}
